\SourceLeaf{251}%
triangulo respondens ad sphaerae superficiem integram ut \(\pm(A+B+C-\pi)\) ad \(4\pi\). Hoc theorema, quod ni fallimur ad elegantissima in theoria superficierum curvarum referendum esse videtur, etiam sequenti modo enuntiari potest:

\textit{Excessus summae angulorum trianguli a lineis brevissimis in superficie curva concavo-concava formati ultra \(180^0\), vel defectus summae angulorum trianguli a lineis brevissimis in superficie curva concavo-convexa formati a \(180^0\) mensuratur per aream partis superficiei sphaericae, quae illi triangulo per directiones normalium respondet, si superficies integra 720 gradibus aequiparatur.}

Generalius in quovis polygono \(n\) laterum, quae singula formantur per lineas brevissimas, excessus summae angulorum supra \(2n-4\) rectos, vel defectus a \(2n-4\) rectis (pro indole curvaturae superficiei), aequatur areae polygoni respondentis in superficie sphaerica, dum tota superficies sphaerae 720 gradibus aequiparatur, uti per discerptionem polygoni in triangula e theoremate praecedenti sponte demanat.
\Article{21}
Restituamus characteribus \(p,q,E,F,G,\omega\) significationes generales, quibus supra accepti fuerant, supponamusque, indolem superficiei curvae praeterea alio simili modo per duas alias variabiles \(p',q'\) determinari, ubi elementum lineare indefinitum exprimatur per
\[
\BareRoot{(E'\mathrm{d}p'^2+2F'\mathrm{d}p'\mathord{.}\mathrm{d}q'+G'\mathrm{d}q'^2)}
\]
Ita cuivis puncto superficiei per valores determinatos variabilium \(p,q\) definito respondebunt valores determinati variabilium \(p',q'\), quocirca hae erunt functiones ipsarum \(p,q\), e quarum differentiatione prodire supponemus
\[
\begin{aligned}
\mathrm{d}p'&=\alpha\mathrm{d}p+\beta\mathrm{d}q\\
\mathrm{d}q'&=\gamma\mathrm{d}p+\delta\mathrm{d}q
\end{aligned}
\]
Iam proponimus nobis investigare significationem geometricam horum coëfficientium \(\alpha,\beta,\gamma,\delta\).

\textit{Quatuor} itaque nunc systemata linearum in superficie curva concipi possunt, pro quibus resp. \(q,p,q',p'\) sint constantes. Si per punctum determinatum, cui respondent variabilium valores \(p,q,p',q'\). quatuor lineas ad singula illa systemata pertinentes ductas supponimus, harum elementa, variationibus positivis
